\begin{abstract}
This paper presents a philosophical framework where everything that exists is made of experience. Reality consists of fundamental experiential units called \emph{vibes}, each feeling either pleasure or pain in each moment of time (\emph{beats}). These vibes connect to each other through \emph{links}, creating vast networks where each vibe influences and is influenced by its neighbors. Physical matter, consciousness, and in theory all phenomena we observe could then be considered potentially to emerge as patterns within this living, experiencing mesh of interconnected vibes. This framework offers a speculative philosophical lens for understanding consciousness, reality, and religious experience. This is explicitly not a scientific theory yet with testable predictions, but rather a exploration that draws inspiration from various analytical traditions while offering its own distinctive perspective on the fundamental nature of existence. The framework invites consideration about how all aspects of reality, from quantum phenomena to human consciousness to spiritual experience, might be understood as all having the same basic architecture.
\end{abstract}
